\subsection{Kiến trúc Pipeline và Logic Điều khiển}
\subsubsection{Thanh ghi đệm (Pipeline Registers) và Control Unit}
% =============================================================================
% SLIDE 1: TỔNG QUAN KIẾN TRÚC MODULAR
% =============================================================================
\begin{frame}{Pipeline Registers \& Control Unit}
    \begin{block}{Tư duy thiết kế: Phân rã cấu trúc (Structural Separation)}
        Tách biệt hoàn toàn giữa \textbf{Datapath} (Đường dữ liệu), \textbf{Control Unit} (Khối điều khiển tổ hợp), và \textbf{Pipeline Registers} (Vách ngăn đồng bộ).
    \end{block}

    \begin{columns}[T]
        \begin{column}{0.5\textwidth}
            \textbf{Control Unit (Thuần tổ hợp):}
            \begin{itemize}
                \item Giải mã mã lệnh 32-bit từ tầng IF.
                \item Sinh tín hiệu điều khiển 1-bit ($reg\_write, mem\_read,\dots$).
                \item \textit{[Bạn sẽ paste code Control Unit tại đây hoặc slide phụ]}
            \end{itemize}
        \end{column}
        
        \begin{column}{0.5\textwidth}
            \textbf{Pipeline Registers (Chốt chặn Clock):}
            \begin{itemize}
                \item 4 khối đệm: \texttt{IF/ID}, \texttt{ID/EX}, \texttt{EX/MEM}, \texttt{MEM/WB}.
                \item Đóng vai trò là vách ngăn phân chia 5 tầng xử lý.
                \item Tích hợp chân \texttt{hold} (Stall) và \texttt{clear} (Flush).
            \end{itemize}
        \end{column}
    \end{columns}
\end{frame}